\documentclass[11pt]{article}

\usepackage[margin=0.9in]{geometry}
\usepackage{amsmath,amssymb,amsthm}
\usepackage{graphicx}
\usepackage{microtype}
\usepackage[hidelinks]{hyperref}
\usepackage{enumitem}
\setlength{\emergencystretch}{2em}

\newtheorem{theorem}{Theorem}
\newtheorem{lemma}{Lemma}
\newcommand{\M}{\mathcal M}
\newcommand{\supp}{\operatorname{supp}}

\title{The Type L hypothesis is unnecessary for\\topological amenability of sub-semihypergroups}
\author{Candidate full solution for arXiv:2211.01479}
\date{19 July 2026}

\begin{document}
\maketitle

\noindent\textbf{Status.}
Candidate full solution, likely valid; human expert review is required.

\section{Source question}

The source is C.~Bandyopadhyay, \emph{Topological Amenability of
Semihypergroups}, arXiv:2211.01479, subsequently published in
\emph{Forum Mathematicum} 36 (2024), no.~4, 943--954 \cite{source}.
The final remark on page 16 of the arXiv PDF asks whether the Type L assumption
in Theorem~3.5 can be removed.

\begin{center}
\includegraphics[width=0.96\textwidth]{figures/open_problem_crop.png}
\end{center}

\noindent\textbf{Transcription of the target sentence.}
``Moreover, although Theorem 3.5 extends the results of [21] even for the case
of locally compact (semitopological) semigroups, it is not known whether in the
setting of a semihypergroup, the restriction that $H$ be a Type L
sub-semihypergroup, can be lifted.''

Here ``lifted'' is read in context as ``removed.''  Theorem~3.5 states that for
a closed sub-\allowbreak semihypergroup $H$ of a semitopological semihypergroup $K$, the
following are equivalent: $H$ has a topological left-invariant mean (TLIM), and
there is a topological left $H$-invariant mean $m$ on $M(K)^*$ with
$m(\chi_H)=1$.  Under the additional Type L hypothesis, the source also proves
equivalence with the formally weaker condition $m(\chi_H)>0$.

\section{Definitions and notation}

We follow the source.  For a locally compact space $X$, $M(X)$ is the Banach
lattice of complex regular Borel measures, $P(X)$ is its set of probability
measures, and convolution is written $*$.  If $H\subseteq K$ is closed,
$\mu_H$ denotes restriction of $\mu\in M(K)$ to $H$, while $\sigma^e$ denotes
extension by zero of $\sigma\in M(H)$ to $K$.  The mass functional of $H$ is
\[
  \chi_H(\mu)=\mu(H) \qquad (\mu\in M(K)).
\]
For $\rho\in M(K)$ and $F\in M(K)^*$,
\[
  (L_\rho F)(\mu)=F(\rho*\mu).
\]
A mean $m$ on $M(K)^*$ is \emph{topological left $H$-invariant} if
$m(L_\rho F)=m(F)$ for every $F\in M(K)^*$ and every $\rho\in P(K)$ with
$\supp\rho\subseteq H$.

Define the positive restriction map
\[
 e_H:M(H)^*\longrightarrow M(K)^*,
 \qquad e_H(\phi)(\mu)=\phi(\mu_H).
\]

\section{Proof Intuition}

The source proof estimates how restriction fails to commute with left
convolution and invokes Type L to rewrite an inverse-set expression.  There is
a more direct description of the error.  If the left factor is supported in
$H$, then the part of the right-hand measure already in $H$ stays in $H$,
because $H*H\subseteq H$.  Thus the entire restriction defect is precisely the
mass that starts in $H^c$ and is transported into $H$.

This incoming mass is measured by the positive functional
$L_{\nu^e}\chi_H-\chi_H$.  The restriction defect is dominated in lattice
modulus by a scalar multiple of that functional.  But an $H$-invariant mean
annihilates it.  Consequently the normalized restriction of any invariant
mean assigning positive mass to $H$ is automatically a TLIM on $H$.

\section{Result}

\begin{theorem}[Theorem 3.5 without Type L]\label{thm:main}
Let $K$ be a semitopological semihypergroup and let $H\subseteq K$ be a closed
sub-semihypergroup.  The following are equivalent:
\begin{enumerate}[label=\textup{(\arabic*)}]
  \item $H$ admits a topological left-invariant mean;
  \item there is a topological left $H$-invariant mean $m\in\M(K)$ such that
        $m(\chi_H)=1$;
  \item there is a topological left $H$-invariant mean $m\in\M(K)$ such that
        $m(\chi_H)>0$.
\end{enumerate}
In particular, the Type L hypothesis in the final part of the source's
Theorem~3.5 is unnecessary.
\end{theorem}

\begin{proof}
The implications $(1)\Rightarrow(2)\Rightarrow(3)$ are already contained in
the source and do not use Type L.  We prove $(3)\Rightarrow(1)$.

Let $m$ be as in (3), set $\alpha=m(\chi_H)>0$, and define
\[
  m_0=\alpha^{-1}m\circ e_H\in M(H)^{**}.
\]
The map $e_H$ is positive.  If $\mathbf 1_H\in M(H)^*$ is the total-mass
functional, then $e_H(\mathbf 1_H)=\chi_H$, and hence
$m_0(\mathbf 1_H)=1$.  Thus $m_0$ is a mean on $M(H)^*$.

Fix $\nu\in P(H)$ and write $\bar\nu=\nu^e\in P(K)$.  Introduce the positive
operator
\[
 Q_\nu:M(K)\longrightarrow M(H),\qquad
 Q_\nu(\mu)=\bigl(\bar\nu*(\mu_{H^c})^e\bigr)_H.
\]
For every $\phi\in M(H)^*$, bilinearity of convolution and
$H*H\subseteq H$ give the exact identity
\begin{equation}\label{eq:defect}
 e_H(L_\nu\phi)-L_{\bar\nu}e_H(\phi)=-\phi\circ Q_\nu.
\end{equation}
Indeed, for $\mu\in M(K)$ we decompose
$\mu=(\mu_H)^e+(\mu_{H^c})^e$.  The measure
$\bar\nu*(\mu_H)^e$ is supported in $H$, and its restriction to $H$ is
$\nu*\mu_H$.  Subtracting this term from
$(\bar\nu*\mu)_H$ leaves exactly $Q_\nu(\mu)$.

Now set
\[
 g_\nu=L_{\bar\nu}\chi_H-\chi_H\in M(K)^*.
\]
For every $\mu\in M^+(K)$, the part $(\mu_H)^e$ stays in $H$ after
convolution by the probability measure $\bar\nu$, with unchanged total mass.
Therefore
\begin{align}\label{eq:incoming}
 g_\nu(\mu)
 &= (\bar\nu*\mu)(H)-\mu(H)\\
 &= \bigl(\bar\nu*(\mu_{H^c})^e\bigr)(H)
  = \|Q_\nu(\mu)\|\geq0.
\end{align}
Thus $g_\nu$ is positive.  Since $Q_\nu$ is positive, the usual variation
inequality gives $|Q_\nu(\mu)|\leq Q_\nu(|\mu|)$.  Hence, in the dual Banach
lattice $M(K)^*$,
\begin{equation}\label{eq:domination}
  |\phi\circ Q_\nu|\leq \|\phi\|g_\nu.
\end{equation}

Because $\supp\bar\nu\subseteq H$, topological left $H$-invariance of $m$
implies
\[
 m(g_\nu)=m(L_{\bar\nu}\chi_H)-m(\chi_H)=0.
\]
Positivity of $m$ and \eqref{eq:domination} now yield
\[
 |m(\phi\circ Q_\nu)|
 \leq m(|\phi\circ Q_\nu|)
 \leq \|\phi\|m(g_\nu)=0.
\]
Applying $m$ to \eqref{eq:defect} and using $H$-invariance once more, we get
\begin{align*}
 m_0(L_\nu\phi)-m_0(\phi)
 &=\alpha^{-1}m\bigl(e_H(L_\nu\phi)-e_H(\phi)\bigr)\\
 &=\alpha^{-1}m\bigl(e_H(L_\nu\phi)-L_{\bar\nu}e_H(\phi)\bigr)\\
 &\quad+\alpha^{-1}m\bigl(L_{\bar\nu}e_H(\phi)-e_H(\phi)\bigr)=0.
\end{align*}
This holds for every $\nu\in P(H)$ and $\phi\in M(H)^*$, so $m_0$ is a TLIM
on $H$.
\end{proof}

\section{Verification, scope, and novelty}

The proof uses only positivity, bilinearity, $H*H\subseteq H$, and topological
$H$-invariance.  It does not use the Type L identity, an involution, an
identity element, commutativity, or any hypergroup axiom.  The crucial point is
the exact defect formula \eqref{eq:defect}; the Banach-lattice estimate then
turns the defect into a positive mean-zero quantity.

This result answers only the Type L question.  It does not address the other
question in the same source remark, namely whether ordinary left amenability
implies topological left amenability for general semitopological
semihypergroups.

A bounded novelty search covered the run's registry and packet indexes, exact
phrases from the source question, the local arXiv corpus, the source's 2024
publication record, the author's later arXiv:2404.18261 and 2024 amenability
survey metadata, and web searches combining ``Type L'', ``sub-semihypergroup'',
and ``topological left $H$-invariant mean''.  No later answer was found.  This
does not establish originality.

\paragraph{Human-review recommendation.}
Review at high priority.  The main checks are the dual Banach-lattice modulus
step in \eqref{eq:domination} and agreement of the source's complex-measure
conventions with the restriction operator $Q_\nu$.  A specialist citation
search should also check whether the same incoming-mass argument appears under
different terminology.

\begin{thebibliography}{9}
\bibitem{source}
C.~Bandyopadhyay,
\emph{Topological Amenability of Semihypergroups},
arXiv:2211.01479; Forum Math. \textbf{36} (2024), no.~4, 943--954.
\end{thebibliography}

\end{document}
